\scanpage{146}
\uline{NB} On a, \deleted{pour} la structure d'Alg. \inserted{\uline{\unsure{commu|comm.|commut.}}} canonique (produit tens.) sur $P=H\otimes H'$, trois sous-Alg quadratiques
\[
 H,\ H',\ H''
\]
\margin{(où $H''=H\wedge H'$ on suppose maintenant $(2,b,b')k=k$)}
\[
 (P/k\simeq H/\deleted{\unsure{k}}k\oplus H'/k\oplus H''/k\quad\text{comme modules})
\]
et trois involutions \unsure{mutuellt|mutuellement} commutantes
\[
 \sigma_H\deleted{\unsure{=}},\ \sigma_{H'}\quad\text{et}\quad\sigma_{H''}
\]
\margin{($\sigma_{H''}$ défini comme $\sigma_H\sigma_{H'}$)}
de telle façon que
\[
 H=P^{\sigma_H},\quad H'=P^{\sigma_{H'}},\quad H''=P^{\sigma_{H''}}.
\]
\deleted{et que $P$ soit un} \inserted{On a.} $H\cap H'=H'\cap H''=H''\cap H=k$ (\unsure{univ|union}$^{\mathrm t}$), \deleted{\unsure{loc. libre|loc. librement} de rang $4$ et} $P$ est \deleted{engendré} comme algèbre par $H$ et $H'$ (d'où $H\otimes H'\overset{\sim}{\longrightarrow}P$), \deleted{et $\unsure{P|B}$ est aussi} (ce qui caractérise $H''$ comme $H\wedge H'$).

Pour que $P$ soit aussi engendré par $H$ et $H''$ \uline{et} $H'$ et $H''$, il f et s. que $H,H'$ soient étales (et alors $H''$ l'est aussi). Si $H$ est étale, $H'$ non, \inserted{alors $H''$ non plus, et} alors $P$ est engendré par $H,H'$ et $H,H''$ (la relation entre $H'$ et $H''$ p.r. à $H$ est symétrique : $H''\simeq H\wedge H'\iff H'\simeq H\wedge H''$) mais \uline{non} par $H'$ et $H''$. En effet prenant une base $u$ de $\deleted{\unsure{H''}}H$, $v$ de $H'$ avec :
\[
 u^2=u,\ v^2=0\qquad(\text{donc }b=1,\ c=b'=c'=0)
\]
et posant
\[
 w=u\ast v=v-2uv
\]
on aura \deleted{bien} \deleted{\unsure{vw} $w^2=0,\ vw=$} base de $P$
\[
 1,u,v,w
\]
mais on a
\[
 v^2=w^2=vw=0
\]
donc $k[\unsure{w|v}]=H'$ et $k[w]=H''$ engendrent ss-alg
% archive sheet number: 145

